Вложенная документная модель MongoDB строится вокруг данных, которые обычно
читаются вместе. Для нее используются четыре коллекции: \texttt{areas},
\texttt{cameras}, \texttt{events} и \texttt{camera\_telemetry}.

В коллекции \texttt{areas} документ территории содержит массив
\texttt{zones}. Зоны вложены в территорию, потому что они не имеют
самостоятельного потока записей и всегда принадлежат одной территории.

Камеры вынесены в коллекцию \texttt{cameras}. Камера не встраивается в зону,
так как она используется не только при описании территории, но и в событиях,
и в телеметрии. Связь камеры с местом установки хранится внутри документа
камеры: в него встроены краткие сведения о территории и зоне. В тот же
документ помещены \texttt{position} и \texttt{settings}, поскольку положение
камеры и настройки обычно читаются вместе с самой камерой.

События хранятся в коллекции \texttt{events}. В документ события встроены
территория, зона, список камер, параметры события \texttt{payload} и массив
обнаруженных объектов \texttt{payload.objects}. Поэтому событие можно читать
сразу вместе с контекстом, не восстанавливая его через отдельные коллекции.

Телеметрия хранится в отдельной коллекции \texttt{camera\_telemetry}, потому
что это быстро растущий временной ряд. Каждая запись телеметрии содержит
краткие сведения о камере, территории и зоне, а технические показатели
помещены во вложенный документ \texttt{metrics}. Такое дублирование увеличивает
объем данных, но уменьшает число обращений при чтении.

\newpage
